\documentclass[11pt]{article}
\usepackage[margin=1in]{geometry}
\usepackage{amsmath,amssymb}
\usepackage[hidelinks]{hyperref}

\newcommand{\co}{\operatorname{co}}

\begin{document}

\section*{Human Review: Initial sign vectors in \texorpdfstring{$c_0$}{c0}}

\noindent Original automatic proof: \texttt{solution\_packet.pdf}

\noindent Checked date: 18 June 2026

\noindent Status: Manually checked.

\noindent Verdict: The proof appears correct.

\section*{Human Comments}

The proof appears correct.  It gives a simple counterexample to the question
in Remark 5.9(iii) of Kurka's paper, under the statement as printed.

The example is
\[
  X=c_0,
\]
and \(W\) is the union of the finite initial sign cubes
\[
  W_n=
  \left\{
  (\varepsilon_1,\ldots,\varepsilon_n,0,0,\ldots):
  \varepsilon_i\in\{-1,1\}
  \right\}.
\]
The construction is very natural once the example has been found: the sets
\(W_n\) are exactly the vertices of the finite cubes in the first \(n\)
coordinates.

\section*{Convex Generation}

The verification of
\[
  \overline{\co} W=B_{c_0}
\]
is correct.  For each \(n\), the convex hull of \(W_n\) is the cube
\[
  \{x\in c_0:\operatorname{supp}x\subseteq\{1,\ldots,n\},
    \ |x_i|\leq 1\}.
\]
Therefore every finitely supported element of \(B_{c_0}\) belongs to
\(\co W\).  Since finitely supported elements are dense in \(B_{c_0}\) by
coordinate truncation, the closed convex hull of \(W\) is the whole unit ball.

\section*{Weakly Convergent Sequences in \texorpdfstring{$W$}{W}}

The weak-sequential part is also correct.  Suppose that
\[
  w_j\in W_{n_j}
\]
is weakly convergent in \(c_0\).  Weak convergence in \(c_0\) implies
coordinatewise convergence, because the coordinate functionals belong to
\(c_0^*=\ell_1\).

If the lengths \(n_j\) were unbounded, one could pass to a subsequence with
\(n_j\to\infty\).  Then for every fixed coordinate \(k\), the \(k\)-th
coordinate of \(w_j\) is eventually always a sign, hence lies in
\(\{-1,1\}\).  Since this coordinate sequence converges, it must eventually be
constant, and the weak limit would have
\[
  |x_k|=1
\]
for every \(k\).  This is impossible for \(x\in c_0\).

Thus the lengths \(n_j\) are bounded.  Consequently the sequence \((w_j)\)
takes values in a finite set.  A finite-valued weakly convergent sequence in a
Hausdorff topology is eventually constant, hence norm convergent.

\section*{Failure of the Schur Property}

Finally, \(c_0\) does not have the Schur property.  The unit vector basis
\((e_n)\) is weakly null in \(c_0\), since for every
\(a=(a_n)\in\ell_1=c_0^*\),
\[
  a(e_n)=a_n\to0.
\]
But \(\|e_n\|_\infty=1\) for all \(n\), so \((e_n)\) is not norm null.

\section*{Review Outcome}

The proof appears correct.  The construction is elementary, and the idea is
clear once one notices the right set \(W\): use initial sign vectors to
generate the unit ball of \(c_0\), while forcing weakly convergent sequences in
\(W\) to have bounded length and hence to be eventually constant.

This should be recorded as an apparently correct counterexample to the
question as printed, subject only to the usual semantic check that no extra
condition on \(W\) was intended.

\end{document}
